\section{Hilbert functions}
\label{am-ch11-hilbert}
\begin{definition}[Poincare series]
\label{am-def-poincare}
\source{atiyah-macdonald-commutative-algebra:page-0125}\uses{}
For a Noetherian graded ring $A=\bigoplus A_n$ and finite graded module
$M=\bigoplus M_n$, and an additive function $\lambda$ on finite $A_0$-modules,
the Poincare series is $P(M,t)=\sum_{n\ge0}\lambda(M_n)t^n$.  Let $d(M)$ be
the order of its pole at $t=1$.
\end{definition}
\begin{theorem}[Hilbert--Serre]
\label{am-thm-11-1}
\dcref{11.1}
\source{atiyah-macdonald-commutative-algebra:page-0126}\uses{}
The Poincare series has the form
$f(t)/\prod_i(1-t^{k_i})$ with $f(t)\in\Z[t]$.
\end{theorem}
\begin{proof}
Induct on the number of homogeneous algebra generators.  Multiplication by the
last generator gives exact sequences on homogeneous pieces; additivity of
$\lambda$ yields a relation expressing $(1-t^{k})P(M,t)$ using the Poincare
series of a kernel and cokernel annihilated by that generator.  Apply induction
to both, with the zero-generator case finite length in degree.
\uses{am-prop-2-11}
\end{proof}
\begin{corollary}[Hilbert polynomial in standard grading]
\label{am-cor-11-2}
\dcref{11.2}
\source{atiyah-macdonald-commutative-algebra:page-0126}\uses{}
If all algebra generators have degree one, then $\lambda(M_n)$ agrees for large
$n$ with a polynomial of degree $d(M)-1$ (degree of the zero polynomial is
$-1$).
\uses{am-thm-11-1}
\end{corollary}
\begin{proof}
Cancel powers of $(1-t)$ so the numerator has nonzero value at $1$, and expand
$(1-t)^{-d}$ using binomial coefficients.  Coefficients are eventually a
polynomial of degree $d-1$ with nonzero leading coefficient.
\end{proof}
\begin{proposition}[Regular homogeneous element]
\label{am-prop-11-3}
\dcref{11.3}
\source{atiyah-macdonald-commutative-algebra:page-0127}\uses{}
If $x\in A_k$ is not a zero-divisor on $M$, then
$d(M/xM)=d(M)-1$.
\end{proposition}
\begin{proof}
In the multiplication exact sequence the kernel vanishes; the Hilbert--Serre
relation therefore divides the pole order by one.
\uses{am-thm-11-1}
\end{proof}
\begin{proposition}[Hilbert polynomial for local filtrations]
\label{am-prop-11-4}
\dcref{11.4}
\source{atiyah-macdonald-commutative-algebra:page-0127}\uses{}
Let $(A,\mathfrak{m})$ be Noetherian local, $\mathfrak{q}$ $\mathfrak{m}$-primary, and $M$ finite with a
stable $\mathfrak{q}$-filtration.  Then $M/M_n$ has finite length; for large $n$ its
length is a polynomial $g(n)$ of degree at most the least number of generators
of $\mathfrak{q}$; and the degree and leading coefficient are independent of the stable
filtration.
\end{proposition}
\begin{proof}
The graded pieces are finite modules over the Artin ring $A/\mathfrak{q}$, so have finite
length.  The associated graded module is finite (Proposition~\ref{am-prop-10-22});
Corollary~\ref{am-cor-11-2} makes the successive lengths polynomial, and
summing gives $g(n)$.  Bounded difference of stable filtrations (Lemma~\ref{am-lem-10-6})
gives inequalities between shifted polynomials, forcing equal degree and leading
coefficient.
\uses{am-prop-10-22,am-cor-11-2,am-lem-10-6,am-thm-8-5}
\end{proof}
\begin{corollary}[Characteristic polynomial]
\label{am-cor-11-5}
\dcref{11.5}
\source{atiyah-macdonald-commutative-algebra:page-0128}\uses{}
For an $\mathfrak{m}$-primary $\mathfrak{q}$, $\length(A/\mathfrak{q}^n)$ is eventually a polynomial
$\chi_\mathfrak{q}(n)$ of degree at most the least number of generators of $\mathfrak{q}$.
\uses{am-prop-11-4}
\end{corollary}
\begin{proof}
Apply Proposition~\ref{am-prop-11-4} to $M=A$ with its $\mathfrak{q}$-adic filtration.
\end{proof}
\begin{proposition}[Independence of the primary ideal]
\label{am-prop-11-6}
\dcref{11.6}
\source{atiyah-macdonald-commutative-algebra:page-0128}\uses{}
For $\mathfrak{m}$-primary $\mathfrak{q}$, $\deg\chi_\mathfrak{q}=\deg\chi_\mathfrak{m}$.
\end{proposition}
\begin{proof}
For some $r$, $\mathfrak{m}\supseteq\mathfrak{q}\supseteq\mathfrak{m}^r$.  Thus
$\chi_\mathfrak{m}(n)\le\chi_\mathfrak{q}(n)\le\chi_\mathfrak{m}(rn)$ for large $n$, forcing equal degrees.
\end{proof}
\begin{definition}[The integer $d(A)$]
\label{am-def-hilbert-dimension}
\source{atiyah-macdonald-commutative-algebra:page-0128}\uses{}
The common degree of the characteristic polynomials $\chi_\mathfrak{q}(n)$, as $\mathfrak{q}$
ranges over the $\mathfrak{m}$-primary ideals, is denoted by $d(A)$.
\uses{am-prop-11-6}
\end{definition}
\begin{proposition}[First dimension inequality]
\label{am-prop-11-7}
\dcref{11.7}
\source{atiyah-macdonald-commutative-algebra:page-0128}\uses{}
For a Noetherian local ring, the least number of generators of an
$\mathfrak{m}$-primary ideal is at most $d(A)$.
\uses{am-cor-11-5,am-prop-11-6}
\end{proposition}
\begin{proof}
Apply Proposition~\ref{am-prop-11-4} to a minimal generating set and compare
the degree of the resulting characteristic polynomial.
\end{proof}

\section{Dimension of Noetherian local rings}
\label{am-ch11-local-dimension}
\begin{proposition}[Quotient by a regular element]
\label{am-prop-11-8}
\dcref{11.8}
\source{atiyah-macdonald-commutative-algebra:page-0129}\uses{}
If $x$ is a non-zero-divisor on a finite module $M$, then
$\deg\chi_{\mathfrak{q}}(M/xM)\le\deg\chi_{\mathfrak{q}}(M)-1$.
\end{proposition}
\begin{proof}
Set $N=xM\simeq M$ and filter $N$ by $N\cap\mathfrak{q}^nM$.  The exact sequence
$0\to N/(N\cap\mathfrak{q}^nM)\to M/\mathfrak{q}^nM\to(M/xM)/\mathfrak{q}^n(M/xM)\to0$ and
Artin--Rees show the first term has the same leading polynomial as $M$;
subtract lengths.
\uses{am-prop-10-9,am-prop-11-4}
\end{proof}
\begin{corollary}[Dimension drops by a regular element]
\label{am-cor-11-9}
\dcref{11.9}
\source{atiyah-macdonald-commutative-algebra:page-0129}\uses{}
If $x$ is a non-zero-divisor in a Noetherian local ring $A$, then
$d(A/(x))\le d(A)-1$.
\uses{am-prop-11-8}
\end{corollary}
\begin{proof}
Take $M=A$ in Proposition~\ref{am-prop-11-8}.
\end{proof}
\begin{proposition}[Dimension bounded by Hilbert degree]
\label{am-prop-11-10}
\dcref{11.10}
\source{atiyah-macdonald-commutative-algebra:page-0129}\uses{}
For a Noetherian local ring, $\dim A\le d(A)$.
\end{proposition}
\begin{proof}
Induct on $d(A)$.  If $d=0$, lengths $\length(A/\mathfrak{m}^n)$ stabilize, so
Nakayama gives $\mathfrak{m}^n=0$ and $A$ is Artin.  For $d>0$, a prime chain
$\mathfrak{p}_0\subsetneq\cdots\subsetneq\mathfrak{p}_r$ and $x\in\mathfrak{p}_1\setminus\mathfrak{p}_0$ reduce to
the domain $A/\mathfrak{p}_0$ and its quotient by $x$.  The regular-element inequality
and induction bound $r-1\le d-1$.
\uses{am-cor-11-9,am-prop-2-6,am-thm-8-5}
\end{proof}
\begin{corollary}[Finite local dimension]
\label{am-cor-11-11}
\dcref{11.11}
\source{atiyah-macdonald-commutative-algebra:page-0129}\uses{}
Every Noetherian local ring has finite dimension.
\uses{am-prop-11-10}
\end{corollary}
\begin{proof}
The Hilbert degree is finite.
\end{proof}
\begin{corollary}[Finite prime height]
\label{am-cor-11-12}
\dcref{11.12}
\source{atiyah-macdonald-commutative-algebra:page-0130}\uses{}
Every prime in a Noetherian ring has finite height; consequently prime ideals
satisfy the descending chain condition.
\uses{am-cor-11-11,am-cor-3-12}
\end{corollary}
\begin{proof}
The height of $\mathfrak{p}$ is $\dim A_\mathfrak{p}$ and is finite by Corollary~\ref{am-cor-11-11}.
\end{proof}
\begin{proposition}[Parameters]
\label{am-prop-11-13}
\dcref{11.13}
\source{atiyah-macdonald-commutative-algebra:page-0130}\uses{}
If $\dim A=d$ for a Noetherian local ring $(A,\mathfrak{m})$, there is an
$\mathfrak{m}$-primary ideal generated by $d$ elements.  Thus $d\le s(A)$, where $s(A)$
is the least number of generators of an $\mathfrak{m}$-primary ideal.
\end{proposition}
\begin{proof}
Choose $x_i\in\mathfrak{m}$ inductively outside the finite union of minimal primes of
height $i-1$ over the previous ideal (prime avoidance).  Every prime over
$(x_1,\ldots,x_i)$ then has height at least $i$.  After $d$ steps the only
prime over the ideal is $\mathfrak{m}$, so it is $\mathfrak{m}$-primary.
\uses{am-prop-1-11,am-cor-11-12}
\end{proof}
\begin{theorem}[Dimension theorem]
\label{am-thm-11-14}
\dcref{11.14}
\source{atiyah-macdonald-commutative-algebra:page-0130}\uses{}
For a Noetherian local ring, the following are equal: the maximum length of a
prime chain; the degree of $\chi_\mathfrak{m}(n)=\length(A/\mathfrak{m}^n)$; and the least number
of generators of an $\mathfrak{m}$-primary ideal.
\uses{am-prop-11-7,am-prop-11-10,am-prop-11-13}
\end{theorem}
\begin{proof}
Propositions~\ref{am-prop-11-7}, \ref{am-prop-11-10}, and
\ref{am-prop-11-13} give the three inequalities in a cycle.
\end{proof}
\begin{corollary}[Embedding dimension bound]
\label{am-cor-11-15}
\dcref{11.15}
\source{atiyah-macdonald-commutative-algebra:page-0130}\uses{}
For a Noetherian local ring, $\dim A\le\dim_{A/\mathfrak{m}}(\mathfrak{m}/\mathfrak{m}^2)$.
\uses{am-prop-2-8,am-prop-11-13}
\end{corollary}
\begin{proof}
A basis of $\mathfrak{m}/\mathfrak{m}^2$ generates $\mathfrak{m}$ by Proposition~\ref{am-prop-2-8}; apply
Proposition~\ref{am-prop-11-13}.
\end{proof}
\begin{corollary}[Height of minimal primes over finitely generated ideals]
\label{am-cor-11-16}
\dcref{11.16}
\source{atiyah-macdonald-commutative-algebra:page-0130}\uses{}
If $A$ is Noetherian and $x_1,\ldots,x_r\in A$, every minimal prime over
$(x_1,\ldots,x_r)$ has height at most $r$.
\uses{am-thm-11-14,am-cor-3-12}
\end{corollary}
\begin{proof}
Localize at a minimal prime; the images generate a primary ideal, and apply the
dimension theorem to that local ring.
\end{proof}
\begin{corollary}[Principal ideal theorem]
\label{am-cor-11-17}
\dcref{11.17}
\source{atiyah-macdonald-commutative-algebra:page-0131}\uses{}
If $x$ is neither a unit nor a zero-divisor in a Noetherian ring, every minimal
prime over $(x)$ has height one.
\uses{am-cor-11-16,am-prop-4-7}
\end{corollary}
\begin{proof}
Height is at most one by Corollary~\ref{am-cor-11-16}; height zero would make
$x$ a zero-divisor by Proposition~\ref{am-prop-4-7}, contradiction.
\end{proof}
\begin{corollary}[Dimension after quotient]
\label{am-cor-11-18}
\dcref{11.18}
\source{atiyah-macdonald-commutative-algebra:page-0131}\uses{}
If $x\in\mathfrak{m}$ is a non-zero-divisor in a Noetherian local ring, then
$\dim A/(x)=\dim A-1$.
\uses{am-cor-11-9,am-thm-11-14}
\end{corollary}
\begin{proof}
The upper bound is Corollary~\ref{am-cor-11-9}; parameters for $A/(x)$ lift
to an $\mathfrak{m}$-primary ideal generated by one additional element in $A$, giving
the reverse bound by Theorem~\ref{am-thm-11-14}.
\end{proof}
\begin{corollary}[Completion preserves dimension]
\label{am-cor-11-19}
\dcref{11.19}
\source{atiyah-macdonald-commutative-algebra:page-0131}\uses{}
The adic completion of a Noetherian local ring has the same dimension.
\uses{am-prop-10-15,am-thm-11-14}
\end{corollary}
\begin{proof}
The quotients by powers of the maximal ideal, hence their eventual length
polynomials, are unchanged by completion.
\end{proof}
\begin{proposition}[Independence of parameters]
\label{am-prop-11-20}
\dcref{11.20}
\source{atiyah-macdonald-commutative-algebra:page-0131}\uses{}
If $x_1,\ldots,x_d$ is a system of parameters and a homogeneous polynomial
$f$ of degree $s$ satisfies $f(x_1,\ldots,x_d)\in\mathfrak{q}^{s+1}$ for
$\mathfrak{q}=(x_1,\ldots,x_d)$, then every coefficient of $f$ lies in $\mathfrak{m}$.
\uses{am-thm-11-14,am-prop-11-3}
\end{proposition}
\begin{proof}
Reduce modulo $\mathfrak{m}$ and map the polynomial ring over $A/\mathfrak{q}$ onto
$\operatorname{gr}_\mathfrak{q}(A)$.  A nonzero reduced homogeneous relation would lower
the pole order by Proposition~\ref{am-prop-11-3}, contradicting the dimension
theorem; hence all coefficients reduce to zero.
\end{proof}
\begin{corollary}[Parameters are algebraically independent]
\label{am-cor-11-21}
\dcref{11.21}
\source{atiyah-macdonald-commutative-algebra:page-0132}\uses{}
If a field $k\subseteq A$ maps isomorphically onto $A/\mathfrak{m}$, a system of
parameters is algebraically independent over $k$.
\uses{am-prop-11-20}
\end{corollary}
\begin{proof}
The lowest homogeneous part of a polynomial relation has coefficients in $k$
and Proposition~\ref{am-prop-11-20} forces them into $\mathfrak{m}$, hence to zero.
\end{proof}

\section{Regular local rings}
\label{am-ch11-regular}
\begin{theorem}[Characterizations of regular local rings]
\label{am-thm-11-22}
\dcref{11.22}
\source{atiyah-macdonald-commutative-algebra:page-0132}\uses{}
For a Noetherian local ring of dimension $d$, the following are equivalent:
$\operatorname{gr}_\mathfrak{m}(A)\simeq k[t_1,\ldots,t_d]$;
$\dim_k\mathfrak{m}/\mathfrak{m}^2=d$; and $\mathfrak{m}$ is generated by $d$ elements.  Such a ring is
regular local.
\end{theorem}
\begin{proof}
The first implication reads degree one.  The second gives generators by
Proposition~\ref{am-prop-2-8}.  If $\mathfrak{m}=(x_1,\ldots,x_d)$, the map
$k[t_1,\ldots,t_d]\to\operatorname{gr}_\mathfrak{m}(A)$ is onto, and Proposition~\ref{am-prop-11-20}
shows its kernel is zero.
\uses{am-prop-2-8,am-prop-11-20}
\end{proof}
\begin{lemma}[Domain from a graded domain]
\label{am-lem-11-23}
\dcref{11.23}
\source{atiyah-macdonald-commutative-algebra:page-0132}\uses{}
If $\bigcap_n\mathfrak{a}^n=0$ and $\operatorname{gr}_\mathfrak{a}(A)$ is a domain, then
$A$ is a domain.
\end{lemma}
\begin{proof}
For nonzero $x,y$, choose lowest filtration degrees with nonzero initial forms.
Their product is nonzero in the graded domain, so $xy\ne0$ in $A$.
\end{proof}
\begin{proposition}[Completion and regularity]
\label{am-prop-11-24}
\dcref{11.24}
\source{atiyah-macdonald-commutative-algebra:page-0133}\uses{}
A Noetherian local ring is regular iff its completion is regular.
\uses{am-prop-10-16,am-thm-10-26,am-cor-11-19,am-prop-10-22,am-thm-11-22}
\end{proposition}
\begin{proof}
Completion is Noetherian local with the same dimension and associated graded
ring by Propositions~\ref{am-prop-10-16}, \ref{am-thm-10-26},
\ref{am-cor-11-19}, and \ref{am-prop-10-22}; apply Theorem~\ref{am-thm-11-22}.
\end{proof}

\section{Transcendence degree}
\label{am-ch11-transcendence}
\begin{theorem}[Local dimension of an affine variety]
\label{am-thm-11-25}
\dcref{11.25}
\source{atiyah-macdonald-commutative-algebra:page-0133}\uses{}
For an irreducible affine variety $V$ over an algebraically closed field, the
local dimension $\dim A(V)_\mathfrak{m}$ is independent of the point and equals
$\trdeg_k\Frac A(V)=\dim V$.
\end{theorem}
\begin{proof}
By normalization there is a polynomial ring
$B=k[x_1,\ldots,x_d]\subseteq A(V)$ over which $A(V)$ is integral, with
$d=\dim V$.  The polynomial ring is integrally closed.  For a maximal ideal
and its contraction, Lemma~\ref{am-lem-11-26} identifies local dimensions; at
the origin of affine $d$-space the dimension is $d$ by Theorem~\ref{am-thm-11-14}.
\uses{am-thm-11-14,am-prop-5-17}
\end{proof}
\begin{lemma}[Dimension under integral extension]
\label{am-lem-11-26}
\dcref{11.26}
\source{atiyah-macdonald-commutative-algebra:page-0134}\uses{}
If $B\subseteq A$ are domains, $B$ integrally closed, $A$ integral over $B$,
and $\mathfrak{m}\in\Max A$ contracts to $\mathfrak{n}$, then $\mathfrak{n}$ is maximal and
$\dim A_\mathfrak{m}=\dim B_\mathfrak{n}$.
\end{lemma}
\begin{proof}
Maximality is Corollary~\ref{am-cor-5-8}.  Incomparability gives one inequality
for chains of primes after contraction, and going down (Theorem~\ref{am-thm-5-16})
 lifts every chain in $B_\mathfrak{n}$ to one in $A_\mathfrak{m}$; hence equality.
\uses{am-cor-5-8,am-cor-5-9,am-thm-5-16}
\end{proof}
\begin{corollary}[Global affine dimension]
\label{am-cor-11-27}
\dcref{11.27}
\source{atiyah-macdonald-commutative-algebra:page-0134}\uses{}
For every maximal ideal of the coordinate ring of an irreducible affine variety,
the local dimension equals the dimension of the variety.
\uses{am-thm-11-25}
\end{corollary}
\begin{proof}
This is the assertion of Theorem~\ref{am-thm-11-25} at each point.
\end{proof}
